Stochastic dual dynamic programming (SDDP) approximates the value function of a multistage problem by cutting planes, and neural SDDP learns those cutting planes across problem instances. I apply both to the Analytic Climate Economy (ACE), an integrated assessment model with DICE functional forms and a closed-form solution, so that every result can be checked exactly. With business-as-usual emissions on an exogenous path, each period of ACE is a small convex program, and the gradient of its value, obtained from the multipliers on the cuts, is the costate recursion of the model's equilibrium conditions: the slopes of the cuts are ACE's costates, and the social cost of carbon (SCC) follows from them. Because ACE's value function is affine, SDDP is exact after two iterations; it reproduces the closed-form SCC to within 0.002\% across a wide range of time preferences, damages and carbon-cycle parameters, in under three seconds per parameter draw. A cut network trained on 128 SDDP solutions predicts the cuts for new parameter values; one forward pass with these cuts, without any SDDP iterations, gives the SCC with a median error of 0.2\% (at most 2.5\%), 8 to 19 times more accurate than a pseudo-state deep equilibrium net on the same test draws. Warm-starting SDDP with the predicted cuts makes its first policy nearly optimal but does not reduce the number of iterations, because SDDP from scratch already needs only one backward pass. I discuss what these results imply for climate-economy models without closed forms, where curved value functions would let neural SDDP pay off, and where non-convexities such as tipping points would break SDDP's guarantees.
